% !PRINTED 474
\weberpage{474}
%<<CONT
zu verifizieren; aber auch das Auffinden der Faktoren gelingt
leicht, wenigstens soweit die Rechnungen bis jetzt fortgesetzt sind,
da man häufig in den späteren Fällen früher gefundene Resultate benutzen kann, und überdies die allgemeine Form der Faktoren kennt. Beispiele werden dies erläutern.

\textbf{§\,128. Die ersten Fälle der Berechnung von $f(\sqrt{-m})$.}

Setzen wir in der kubischen Gleichung:
\weq{1}{u^3 - \gamma_2(\omega)\,u - 16 = 0,}
deren Wurzeln nach § 54
\[
f(\omega)^8, \quad -f_1(\omega)^8, \quad -f_2(\omega)^8
\]
sind, $\omega = i$, also [§\,125, (16)], $\gamma_2 = 12$, so folgt:
\[
u^3 - 12u - 16 = 0,
\]
eine Gleichung, deren einzige positive Wurzel $u = 4$ ist, so daß
wir in Übereinstimmung mit § 126 erhalten:
\weq{2}{f(i) = \sqrt[4]{2};}
setzen wir $\omega = \sqrt{-2}$, also $\gamma_2 = 20$ [§\,125, (17)], so erhalten
wir aus (1) die Gleichung:
\[
u^3 - 20u - 16 = 0
\]
mit der einzigen rationalen Wurzel $-4$. Da aber nach den
Sätzen der beiden vorhergehenden Paragraphen $f_1(\sqrt{-2})^8$ rational
sein muß, so ist:
\weq{3}{f_1(\sqrt{-2}) = \sqrt[4]{2}.}

Wir setzen ferner $\omega = e^{\frac{2\pi i}{3}}$, also $\gamma_2(\omega) = 0$ (§ 117) und
erhalten aus (1):
\[
u^3 = 16.
\]
Dieser Gleichung genügt [§\,34, (19)]:
\[
u = -f_2\left(\frac{-1+\sqrt{-3}}{2}\right)^8 = \frac{16\,e^{\frac{2\pi i}{3}}}{f(\sqrt{-3})^8},
\]
woraus der reelle positive Wert:
\weq{4}{f(\sqrt{-3}) = \sqrt[3]{2}.}

Um die übrigen Resultate des § 125 anwenden zu können,
setzen wir
\[
m = 7, \quad 11, \quad 19, \quad 43, \quad 67, \quad 163,
\]
%CONT>>
